\section*{序論}
10月10日朝までの24時間は、大手モデルの新発表よりも、生成AIの事業規模や研究成果をめぐる解釈がX上で目立った。OpenAIの年換算売上を約500億ドルとする報道と従来の約700億ドルという数字の違いが議論され、AI生成数学原稿への批判も拡散した。いずれも一次資料の未確認部分が大きく、投稿上の主張と確定情報を分けて読む必要がある。

企業・研究者からは、小型モデルや文書処理、合成音声検出、LLMの不確実性評価、推論オークション、軌道最適化などの投稿があった。エージェントの実行基盤・評価・決済も話題となった一方、Haiku 5.5や企業向けGemini agentなどは発表本体が観測窓外で、当日は利用者による比較・解説が中心だった。Grokの収集範囲はPartialであり、以下は窓内の投稿を整理した記録である。
